% main.tex  (FP3)
% Fixed Points III: Disturbance--Damping Balance and Stability
% Uses: series.sty (v3) and refs.bib
\documentclass[11pt]{article}
\usepackage{series}

% ---------------------------
% Paper metadata (override series defaults)
% ---------------------------
\renewcommand{\PartRoman}{III}
\hypersetup{
  unicode=true,
  pdftitle={Fixed Points III: Disturbance--Damping Balance and Stability},
  pdfauthor={Peter Nero}
}

% ---------------------------
% Local notation
% ---------------------------
\newcommand{\Aop}{A}

\newcommand{\HoneF}{H^{1}_{F}}
\newcommand{\Ltwo}{L^{2}}


\newcommand{\E}{\mathbb{E}}
\newcommand{\Var}{\mathrm{Var}}

% Bookmark-safe helper (series.sty defines it; keep fallback)
\providecommand{\FPpdf}[2]{\texorpdfstring{#1}{#2}}

\title{Fixed Points III: Disturbance--Damping Balance and Stability}
\author{Peter Nero}
\date{September 2026\\Version 6}

\begin{document}
\maketitle
\seriestagline

\begin{abstract}
We develop disturbance--damping stability on the joint spectral decomposition
of the FP--II internal operators. For each noncoherent joint mode
$\alpha$, the one-sided nonlinear damping margin is
$\gamma_\alpha=d_\alpha-L_\alpha$. Deterministic force amplitude
$f_\alpha$ and stochastic noise power $q_\alpha$ are distinct quantities. We
prove the deterministic input-to-state floor
$\limsup|a_\alpha|\le f_\alpha/\gamma_\alpha$ and the stochastic
second-moment floor
$\limsup\mathbb E|a_\alpha|^2\le q_\alpha/(2\gamma_\alpha)$ when
$\gamma_\alpha>0$. Gaussian invariant laws and necessity of the sign condition
are asserted only for exact Ornstein--Uhlenbeck dynamics or robust worst-case
stability. Bundlewise results apply to $Q\Psi$ and use separate stochastic trace
and deterministic weighted-series conditions. Finally, deterministic
homogenization is made conditional on an enhanced functional central-limit
theorem, tightness, and rough-path convergence; the symmetrized Green--Kubo tensor
is $D=\int_0^\infty(R+R^\ast)\,ds$.
\end{abstract}

\section*{Version 6 Revision Note}
\noindent\textbf{Supersedes.} Version 5, which remains the released edition.
\textbf{Reason and resolution.} Separates the standalone analytic theorems from the physical source-selection problem.
\textbf{Retained results.} The previous mathematical results and proofs are unchanged.
\textbf{Boundary.} This is a contextual clarification in an unreleased revision,
not a new physical-source theorem.

\section*{Revision note for version 5}
\addcontentsline{toc}{section}{Revision note for version 5}
\begin{description}
\item[Supersedes.] \emph{Fixed Points III: Disturbance--Damping Balance and
Stability}, version~4.
\item[Reason.] The corrected deterministic and stochastic estimates in
version~4 were sound, but their parallel notation could still obscure that
force amplitude, noise power, invariant law, and homogenized noise are
different objects.
\item[Resolution.] Version~5 adds an explicit argument map and worked
interpretive discussion of the two disturbance lanes, their summability
requirements, and the separate homogenization gate. It preserves theorem
ownership and separates the public abstract from revision history and
reproducibility metadata.
\item[Retained result.] The deterministic input-to-state floor, stochastic
second-moment floor, OU restrictions, and enhanced invariance-principle
conditions are unchanged.
\item[Remaining boundary.] Coherent-sector disturbances, general nonlinear
invariant measures, and a physical stochastic source law remain unproved.
\end{description}

\section*{Revision note for version 4}
\addcontentsline{toc}{section}{Revision note for version 4}
\begin{description}
\item[Supersedes.] \emph{Fixed Points III: Disturbance--Damping Balance and
Stability}, version~3.
\item[Reason.] Deterministic amplitudes and stochastic powers were conflated,
nonlinear systems were assigned Gaussian/OU conclusions, and the Green--Kubo
normalization and homogenization hypotheses were incomplete.
\item[Resolution.] Version~4 derives separate deterministic and stochastic
floors on the joint modal spectrum, restricts Gaussian claims to exact OU
dynamics, and states the enhanced invariance-principle and rough-path gates.
\item[Retained result.] Positive net modal damping still gives quantitative
noncoherent stability and disturbance floors.
\item[Remaining boundary.] Coherent disturbances, nonlinear invariant laws,
and physical stochastic interpretation require additional assumptions.
\end{description}

\section*{How to read this paper}
\addcontentsline{toc}{section}{How to read this paper}
FP--II separates the coherent sector $P\Psi$ from its damped complement
$Q\Psi$.  This paper asks what happens to the complement when it is continually
disturbed.  The answer depends on what kind of disturbance is meant, so the
paper deliberately develops deterministic forcing and stochastic forcing in
parallel without identifying their parameters.

\paragraph{The central picture in plain language.}
For each joint internal mode, damping removes amplitude at a rate
$d_\alpha$, while nonlinear feedback can return at most
$L_\alpha$ in the one-sided energy estimate.  Their difference
$\gamma_\alpha=d_\alpha-L_\alpha$ is the available margin.  A bounded
deterministic push produces an amplitude floor proportional to
$f_\alpha/\gamma_\alpha$.  Brownian forcing deposits variance rather than a
fixed amplitude and produces a second-moment floor proportional to
$q_\alpha/(2\gamma_\alpha)$.  The similar-looking denominators do not make
$f_\alpha$ and $q_\alpha$ the same quantity.

\paragraph{Argument map.}
Sections~1--2 inherit the joint projector and build a non-double-counted
spectral index for the overlapping vertical operators.  Sections~3--4 derive
the deterministic and stochastic modewise bounds.  Section~5 identifies the
limited settings in which positivity of the margin is also necessary and a
Gaussian invariant law follows.  Section~6 lifts individual mode bounds to
the bundle by summability or covariance-trace conditions.  Section~7 explains
when rapidly mixing deterministic complement modes can instead appear as
effective stochastic forcing of coherent variables.  Sections~8--9 separate
invariant measures from deterministic fixed points and turn the hypotheses
into an execution checklist.

\paragraph{Scope boundary.}
Every modal result in the main stability chain concerns $Q\Psi$.  It neither
contracts coherent modes nor proves that a physical environment supplies
Brownian noise.  The homogenized diffusion is an emergent limit only after
the enhanced functional-CLT, tightness, and rough-path hypotheses have been
verified for the selected dynamics.

A physical q79 application must independently select the damping generator,
coherent projector, invariant domain, gap data, and forcing or noise law.  The
analytic estimates below consume those typed inputs; they do not emit them
from MTT geometry.

\tableofcontents

\section{Scope and inherited framework}

We use the corrected FP--I/II control framework. The compact internal space
carries strongly commuting nonnegative self-adjoint operators $A_1,A_2,A_3$
on one common Hilbert space. Their joint harmonic projector is $P$, with
$Q=I-P$. All stability results below concern $Q\Psi$ unless coherent forcing
is explicitly introduced. The stabilization parameter is not physical
Lorentzian time.

Deterministic fixed points of a time-step map and stochastic invariant measures
are different objects. Existence of the former is inherited from FP--I/II under
their invariant-set and compactness/condensing hypotheses. Existence of the
latter requires a Markov/Feller and tightness argument stated separately below.
For analytic-semigroup and dissipative-flow background, see
\cite{Amann1995,Henry1981,Temam1997}; for stochastic evolution and exact
Ornstein--Uhlenbeck processes, see
\cite{DaPratoZabczyk1992,Gardiner2009}. The companion analytic framework is
\cite{Nero2025FP1}.

\section{Joint modal decomposition}

Strong commutation supplies a joint spectral calculus. In the compact pure-
point case, choose a joint orthonormal basis $\{e_\alpha\}$ such that
\[
A_ne_\alpha=\lambda_\alpha^{(n)}e_\alpha,
\qquad n=1,2,3.
\]
The multi-index $\alpha$ includes all three spectral labels and any
multiplicity label. Define
\[
d_\alpha:=\sum_{n=1}^3\kappa_n\lambda_\alpha^{(n)},
\qquad
\Lambda_\alpha:=\sum_{n=1}^3\lambda_\alpha^{(n)}.
\]
Let $\mathcal I_Q$ contain the indices for which at least one
$\lambda_\alpha^{(n)}$ is positive. Then
\[
Q\Psi(t)=\sum_{\alpha\in\mathcal I_Q}a_\alpha(t)e_\alpha.
\]
This joint index prevents multiple counting when the vertical structures
overlap. It also applies to a genuinely nested realization after its inclusion
maps have been supplied; rank labels alone do not establish nesting.

\begin{assumption}[One-sided modal remainder bound]\label{ass:one-sided}
On the invariant set, the nonlinear modal remainder satisfies
\[
\operatorname{Re}(\overline{a_\alpha}R_\alpha(a))
\le L_\alpha|a_\alpha|^2,
\qquad
\gamma_\alpha:=d_\alpha-L_\alpha.
\]
The remainder may couple modes; the displayed one-sided estimate, rather than
an entrywise linearization identity, is the hypothesis used in the energy
inequalities.
\end{assumption}

\section{Two distinct disturbance classes}

\subsection{Deterministic forcing}

For deterministic input $b_\alpha(t)$ assume
\[
\|b_\alpha\|_{L^\infty(0,\infty)}\le f_\alpha,
\]
where $f_\alpha$ has units of amplitude per stabilization time. The mode
equation is
\begin{equation}
\dot a_\alpha=-d_\alpha a_\alpha+R_\alpha(a)+b_\alpha(t).
\label{eq:det-mode}
\end{equation}

\subsection{Stochastic forcing}

For the scalar real stochastic theorem let $W_\alpha$ be standard Brownian
motions and let $q_\alpha\ge0$ be noise powers. The mode equation is
\begin{equation}
da_\alpha=(-d_\alpha a_\alpha+R_\alpha(a))\,dt
+\sqrt{q_\alpha}\,dW_\alpha(t).
\label{eq:stoch-mode}
\end{equation}
Correlated or infinite-dimensional noise is described by a positive covariance
operator $Q_\xi$; then the scalar $q_\alpha$ are replaced by its matrix entries
and the bundlewise criterion is a weighted trace condition.

The deterministic amplitude $f_\alpha$ and stochastic power $q_\alpha$ are not
interchangeable and are not represented by one common parameter.

\section{Modewise stability}

\begin{theorem}[Deterministic input-to-state bound]\label{thm:det-iss}
Assume Assumption~\ref{ass:one-sided} and
$\gamma_\alpha>0$. Every solution of Equation~\eqref{eq:det-mode} obeys
\[
|a_\alpha(t)|
\le e^{-\gamma_\alpha t}|a_\alpha(0)|
+\frac{f_\alpha}{\gamma_\alpha}
(1-e^{-\gamma_\alpha t}),
\]
and hence
\[
\limsup_{t\to\infty}|a_\alpha(t)|
\le\frac{f_\alpha}{\gamma_\alpha}.
\]
\end{theorem}

\begin{proof}
The upper Dini derivative satisfies
$D^+|a_\alpha|\le-\gamma_\alpha|a_\alpha|+f_\alpha$.
Comparison with the scalar affine equation gives the result.
\end{proof}

\begin{theorem}[Stochastic second-moment bound]\label{thm:stoch-moment}
Assume Assumption~\ref{ass:one-sided} and
$\gamma_\alpha>0$. Every solution of Equation~\eqref{eq:stoch-mode} satisfies
\[
\frac{d}{dt}\mathbb E|a_\alpha|^2
\le-2\gamma_\alpha\mathbb E|a_\alpha|^2+q_\alpha,
\]
and therefore
\[
\mathbb E|a_\alpha(t)|^2
\le e^{-2\gamma_\alpha t}\mathbb E|a_\alpha(0)|^2
+\frac{q_\alpha}{2\gamma_\alpha}(1-e^{-2\gamma_\alpha t}).
\]
\end{theorem}

\begin{proof}
Apply It\^o's formula to $|a_\alpha|^2$, use the one-sided remainder bound,
and take expectations. The martingale term has zero expectation and the
quadratic variation contributes $q_\alpha dt$.
\end{proof}

These two theorems are sufficient stability bounds for nonlinear dynamics.
They do not assert that a nonlinear invariant law is Gaussian or that
$\gamma_\alpha>0$ is necessary for every particular forcing history.

\paragraph{What the two floors mean.}
Both estimates say that positive damping prevents indefinite accumulation in
one mode, but they answer different questions.  The deterministic theorem
bounds every trajectory against a worst-case input amplitude.  The stochastic
theorem bounds an expectation after quadratic variation has injected power.
Neither floor is automatically attained, and neither is a prediction until
the mode, margin, and disturbance data have been obtained from the concrete
operator model.

\section{Exact OU theorem and robust necessity}

\begin{theorem}[Exact scalar OU classification]\label{thm:OU}
For
\[
da=-\gamma a\,dt+\sqrt q\,dW_t,
\qquad q>0,
\]
there is a unique invariant probability law with finite variance if and only if
$\gamma>0$. It is Gaussian with mean zero and variance
\[
\operatorname{Var}(a)=\frac{q}{2\gamma}.
\]
For $\gamma=0$ the variance grows linearly, and for $\gamma<0$ it grows
exponentially.
\end{theorem}

\begin{theorem}[Worst-case deterministic sign criterion]\label{thm:robust}
The scalar family
$\dot a=-da+R(a)+b(t)$ with
$\operatorname{Re}(\bar aR(a))\le L|a|^2$ is uniformly input-to-state stable
for every bounded input and every admissible remainder only if
$\gamma=d-L>0$. For a specified forcing/remainder pair,
$\gamma\le0$ does not by itself prove divergence because cancellation or a
vanishing input may occur.
\end{theorem}

Thus ``if and only if'' belongs to exact OU dynamics or a declared robust
worst-case class, not to arbitrary nonlinear stochastic equations.

\section{Bundlewise noncoherent stability}

Under bounded geometry and the joint spectral calculus, assume
\begin{equation}
c_1\sum_{\alpha\in\mathcal I_Q}(1+\Lambda_\alpha)|a_\alpha|^2
\le\|Q\Psi\|_{H_F^1}^2
\le c_2\sum_{\alpha\in\mathcal I_Q}(1+\Lambda_\alpha)|a_\alpha|^2.
\label{eq:joint-sobolev}
\end{equation}

\subsection{Stochastic trace condition}

For independent exact OU modes define
\[
\Sigma_q:=\sum_{\alpha\in\mathcal I_Q}
(1+\Lambda_\alpha)\frac{q_\alpha}{2\gamma_\alpha}.
\]
If every $\gamma_\alpha>0$ and $\Sigma_q<\infty$, then
\[
\limsup_{t\to\infty}\mathbb E\|Q\Psi(t)\|_{H_F^1}^2
\le c_2\Sigma_q.
\]
For the independent exact OU product this condition is also necessary for a
finite-$H_F^1$ second moment. For nonlinear dynamics it remains a sufficient
bound unless additional lower estimates establish necessity. With correlated
noise the correct replacement is the corresponding weighted covariance trace.

\subsection{Deterministic weighted-series condition}

Define
\[
\Sigma_f:=\sum_{\alpha\in\mathcal I_Q}
(1+\Lambda_\alpha)\frac{f_\alpha^2}{\gamma_\alpha^2}.
\]
If every $\gamma_\alpha>0$ and $\Sigma_f<\infty$, then
\[
\limsup_{t\to\infty}\|Q\Psi(t)\|_{H_F^1}^2
\le c_2\Sigma_f.
\]
This is an amplitude-series condition, not a stochastic trace.

\subsection{Weyl-law check}

If the joint counting function satisfies
$N(\Lambda)\lesssim\Lambda^{d_{\rm eff}/2}$ and
$\gamma_\alpha\gtrsim1+\Lambda_\alpha$, then
$q_\alpha\lesssim(1+\Lambda_\alpha)^{-p}$ with
$p>d_{\rm eff}/2$ is sufficient for $\Sigma_q<\infty$. The exact exponent must
be recomputed if the damping growth or joint multiplicities differ.

No statement in this section controls a disturbance acting directly in
$P\Psi$. Such coherent forcing needs a separate base/coherent stability
theorem.

\section{Deterministic homogenization of coherent dynamics}

This section addresses a different route to stochastic behavior.  The
underlying system can remain deterministic while rapidly mixing
noncoherent modes drive slow coherent variables.  Under a sufficiently strong
limit theorem, the accumulated fast forcing converges to Brownian transport.
Ordinary mixing is not enough: the second iterated integrals determine whether
the limiting equation is the stated Stratonovich equation or carries an
additional bracket drift.

Write $\Psi=X+Y$ with $X=P\Psi$ and $Y=Q\Psi$. For frozen $x$, let the fast
flow for $Y$ have invariant measure $\mu_x$, and decompose
\[
g(x,y)=\bar g(x)+G(x,y),
\qquad
\bar g(x)=\int g(x,y)\,d\mu_x(y).
\]
The correctly scaled slow--fast system is
\begin{align}
\dot X_\varepsilon
&=f(X_\varepsilon)+\bar g(X_\varepsilon)
+\varepsilon^{-1/2}G(X_\varepsilon,Y_\varepsilon),\\
\dot Y_\varepsilon
&=\varepsilon^{-1}L_{X_\varepsilon}(Y_\varepsilon).
\label{eq:fastslow}
\end{align}

For fixed $x$ define the stationary covariance
\[
R_x(s)=\int G(x,Y_s)\otimes G(x,Y_0)\,d\mu_x(Y_0).
\]

\begin{assumption}[Enhanced invariance principle]\label{ass:efclt}
On the selected compact slow domain:
\begin{enumerate}[label=\textup{(H\arabic*)}]
\item the frozen fast flow has a unique invariant measure and sufficient
mixing for an integrable covariance;
\item the centered additive functionals satisfy a functional CLT uniformly in
$x$, and their laws are tight in the required path topology;
\item the lifted first and second iterated integrals converge to a geometric
Brownian rough path, with zero area anomaly for the canonical Stratonovich
statement below;
\item the Poisson equation/corrector and its $x$-derivatives have the bounds
needed to control slow dependence; and
\item initial layers and exits from the compact slow domain are controlled on
$[0,T]$.
\end{enumerate}
For infinite-dimensional coherent spaces, the corresponding Hilbert-space
tightness, trace, and rough-path assumptions must be supplied separately.
\end{assumption}

\begin{theorem}[Conditional homogenized limit]\label{thm:homogen}
Under Assumption~\ref{ass:efclt}, $X_\varepsilon$ converges in law on
$C([0,T])$ to
\[
dX_t=(f(X_t)+\bar g(X_t))\,dt+\sigma(X_t)\circ dW_t,
\]
where
\begin{equation}
D(x):=\sigma(x)\sigma(x)^\ast
=\int_0^\infty\bigl(R_x(s)+R_x(s)^\ast\bigr)\,ds.
\label{eq:green-kubo}
\end{equation}
If the enhanced limit has a nonzero area anomaly, an additional deterministic
bracket drift must be included; if only an ordinary functional CLT is known,
the Stratonovich conclusion is not established.
\end{theorem}

\section{Fixed points versus invariant measures}

For autonomous or $\tau$-periodic deterministic forcing, the corrected
FP--I/II Schauder or Darbo theorem applies after a deterministic invariant set
is proved; the result is respectively a step-fixed or periodic point, with
equilibrium promotion only under the FP--II Lyapunov condition.

For stochastic forcing, the solution defines a Markov semigroup. An invariant
probability measure requires, for example, the Feller property, a Lyapunov
moment bound, and tight Krylov--Bogoliubov averages. Uniqueness requires an
additional irreducibility/coupling or contractivity argument. A stochastic
invariant measure is not called a deterministic fixed point.

\section{Examples and execution checklist}

For a circle operator with eigenvalues $m^2$ and linear damping
$\gamma_m=m^2-L$, exact OU noise gives
$\operatorname{Var}(a_m)=q_m/[2(m^2-L)]$, whereas bounded deterministic forcing
gives $\limsup|a_m|\le f_m/(m^2-L)$. These are different floors.

For a compact Heisenberg nilmanifold in a noncollapsing left-invariant metric
class, the first positive eigenvalue has a uniform lower bound. This supplies
one ingredient in $d_\alpha$ but does not replace the joint-mode or summability
checks.

For every concrete geometry:
\begin{enumerate}
\item construct the joint spectral index and multiplicities;
\item compute $d_\alpha$ and prove the one-sided constants $L_\alpha$;
\item list deterministic amplitudes $f_\alpha$ and stochastic covariance data
$q_\alpha$ or $Q_\xi$ separately;
\item check $\gamma_\alpha>0$ and the appropriate $\Sigma_f$ or covariance
trace;
\item control coherent forcing independently; and
\item for homogenization, verify the enhanced invariance principle rather than
only quoting mixing.
\end{enumerate}

\section{Conclusion}

The damping-margin principle survives, but with precise scope. Positive
$\gamma_\alpha$ gives deterministic input-to-state and stochastic
second-moment bounds. Exact Gaussian invariant laws and sign necessity belong
to exact OU or robust worst-case formulations. Bundlewise results concern
$Q\Psi$ and require either a stochastic trace or deterministic weighted
series. The homogenized Stratonovich equation is conditional on enhanced
functional-CLT/rough-path data and uses the corrected Green--Kubo
normalization. These distinctions prevent stochastic invariant measures from
being conflated with deterministic fixed points.

The paper therefore contributes a translation layer between FP--II spectral
damping and later statistical descriptions.  Its robust content is the pair
of modewise inequalities and their bundlewise summability conditions.  Exact
Gaussianity belongs only to the OU specialization, while emergent diffusion
belongs only to the verified homogenization limit.  Coherent disturbances,
nonlinear invariant-law uniqueness, and a physical identification of the
noise source remain separate model-specific tasks.

\appendix

\section{Exact OU variance computation}

For $da=-\gamma a\,dt+\sqrt q\,dW_t$,
\[
a(t)=e^{-\gamma t}a(0)+\sqrt q\int_0^t e^{-\gamma(t-s)}\,dW_s.
\]
It\^o isometry gives
\[
\mathbb E\left|\sqrt q\int_0^t e^{-\gamma(t-s)}\,dW_s\right|^2
=\frac{q}{2\gamma}(1-e^{-2\gamma t})
\]
when $\gamma>0$.

\section{Uniform spectral gap for noncollapsing nil fibers}\label{app:nilgap}

\begin{proposition}
Let $F=\Gamma\backslash\mathrm{Nil}_3$ be fixed and let
$\mathcal G_\epsilon$ be a compact class of left-invariant metrics satisfying
$\epsilon I\le G\le\epsilon^{-1}I$. Then
\[
\inf_{g\in\mathcal G_\epsilon}\lambda_1(F,g)>0.
\]
The same bound is uniform for a smoothly base-dependent family remaining in
$\mathcal G_\epsilon$.
\end{proposition}

\begin{proof}
The first positive eigenvalue varies continuously on this compact metric
family and is positive for every compact connected fiber, so it attains a
positive minimum; compare the spectral-geometric background in
\cite{Chavel1984,Buser1982,Kato1976,Hebey2000}.
\end{proof}

\bibliographystyle{plain}
\bibliography{refs}

% BEGIN MTT MANAGED COMPUTATIONAL EVIDENCE
\section*{Computational Evidence and Reproducibility}
The numerical and machine-verifiable claims used by this paper are archived in the curated repository, \texttt{https://github.com/PeterNero/mtt-results-repro}. The mapped authority/result identifiers are \texttt{no result rows mapped}. Claim tiers in that capsule distinguish exact derivation, certified numerics, profile replay, conditional results, and open obligations.
% END MTT MANAGED COMPUTATIONAL EVIDENCE

\end{document}
